\section{有限位数楕円境界の中心商に必要な正則性の証明}
\label{CSBGsec:pair-finite-order-background}
この節では、固定した AS 正則な周囲の中心三次商について、
tails 大域次元が3となる係数範囲を構成する。
二つの中心 chart の正則性を楕円境界に沿って貼り合わせ、
有限位数の場合は Rains の中心持上げと有限中心代数の Bertini 定理から非空開集合を得る。
以下で$B$は標数零の体上の左右noetherian connected graded AS正則代数、
$W\in Z(B)$は正次数の正則斉次元、$A=B/(W)$とする。
$\mathcal P_A=\operatorname{thick}\{\pi A(n):n\in\mathbb Z\}$、
$\mathcal Q_A=\mathrm D^b(\qgr A)/\mathcal P_A$と置く。
有限graded自由行列因子化と次数付き特異圏の比較には
\cite[arXiv v1, Theorems~4.7, 5.5]{CSBGCCKM}を用いる。
右加群には反対環を適用し、中心元のnormalizing automorphismは恒等である。
従って有限syzygy後のcoherent対象を二周期因子化で扱える。
以下の補題の引用入力と仮定を保ち、一般の非正則底へ拡張しない。
\begin{lemma}[有限射影因子化の縮約と分母払い]
\label{CSBGlem:gluing-contraction}
$2$が可逆な左右noetherian環$S$と中心正則元$w$を取り、
有限生成射影加群の行列因子化を$d^2=w$と書く。
$S/(w)$上のcokernelが有限射影次元なら、この因子化は縮約可能である。
さらに中心正則元$u$で局所化した後に縮約可能なら、
ある$N$について$u^N\mathrm{id}=dK+Kd$となる。
次数付きの場合も、因子化の次数shiftを含めて同じ主張が成立する。
\end{lemma}
\begin{proof}
$\bar S=S/(w)$上の二周期分解では、cokernelの有限射影次元により
十分高い syzygy が射影的である。周期性から全ての周期 syzygy が射影的となる。
従って分解は分裂し、
$\bar S$上で$d\bar H+\bar Hd=1$となる奇写像を取れる。
有限射影性からこれを$S$上の$H$へ持ち上げると
$dH+Hd=1+wT$となる。
左辺は$d$と可換なので、$w$の正則性により$[d,T]=0$。
ゆえに$H'=H-\frac12dT$は$dH'+H'd=1$を満たす。
局所化後の$H'$の有限個の成分の分母を$u^N$で払い、
$S\hookrightarrow S[u^{-1}]$の単射性を使えば最後の等式を得る。
\end{proof}

\begin{lemma}[有限大域次元からtwistsによる生成へ]
\label{CSBGlem:global-finite-to-perfect}
$A$は次数1生成で、$\operatorname{qgr}A$が有限大域次元を持つなら、
$\mathrm D^b(\operatorname{qgr}A)=\mathcal P_A$である。
\end{lemma}
\begin{proof}
大域次元を$e$とし、graded自由表示を繰り返して
$0\to K\to L_e\to\cdots\to L_0\to M\to0$を取る。
各$L_i$は$\pi A(n)$の有限直和である。
最後のYoneda拡大は$\operatorname{Ext}^{e+1}(M,K)=0$なので、
$M$は有限複体$[L_e\to\cdots\to L_0]$の直和因子である。
line object自体が射影的であるとは仮定していない。
\end{proof}

\begin{proposition}[中心零化元による十分条件]
\label{CSBGprop:global-annihilators}
$u,v\in Z(B)_1$が$A$上正則列をなし、
$\mathrm D^b(\operatorname{qgr}A/(u,v))$が構造対象のtwistsで厚く生成されるとする。
正次数の斉次中心元$p_1,\ldots,p_r\in Z(A)$が
次数付き特異圏の全対象を零化し、ある$N$について
$(u,v)^N\subset(p_1,\ldots,p_r)A$なら、$\mathcal Q_A=0$である。
\end{proposition}
\begin{proof}
graded lift $M$のKoszul複体$K(p_1,\ldots,p_r;M)$のcohomologyは
$(p_1,\ldots,p_r)$に殺され、従って$(u,v)$の冪に殺される。
有限filtrationの各商は$A/(u,v)$加群である。
正則列のKoszul分解から$\pi(A/(u,v))(n)\in\mathcal P_A$なので、
境界の生成性によりこのKoszul複体は$\mathcal P_A$に入る。
一方、$\mathrm D^b(\operatorname{gr}A)\to\mathcal Q_A$は
次数付きperfect複体を殺すから次数付き特異圏を経由する。
従って各$p_i$は$\mathcal Q_A$で零作用し、反復coneは
$M$のshiftとtwistの直和へ分裂する。ゆえに$M$自身も零となる。
\end{proof}

\begin{proposition}[二つの中心chartの貼合せ]
\label{CSBGprop:global-two-charts}
$u,v\in Z(B)_1$が$B$上正則で$A$上正則列をなし、
$A[u^{-1}]_0,A[v^{-1}]_0$が有限大域次元を持つとする。
$\operatorname{qgr}A/(u,v)$の有界導来圏が構造対象のtwistsで厚く生成されるなら、
$\mathcal Q_A=0$である。
\end{proposition}
\begin{proof}
各 coherent 対象の graded lift に有限回の syzygy を取ると、
行列因子化の cokernel を得る。途中の有限 graded 自由部分は
$\mathcal Q_A$ で消えるので、元の対象はこの cokernel の shift と同型である。
中心一次元を逆にした次数零部で因子化を考えると、
chart の有限大域次元により cokernel の射影次元は有限となり、因子化は縮約する。
補題\ref{CSBGlem:gluing-contraction}の分母払いから、対象ごとに選んだ
$u^r,v^s$が$\mathcal Q_A$で零作用する。
二重Koszul coneのcohomologyは$(u^r,v^s)$に殺され、
$(u,v)^{r+s-1}\subset(u^r,v^s)$から境界の有限filtrationを持つ。
前命題と同じ議論でconeはperfectである。
中心作用は自然変換なので、最初の零coneの直和分解とも可換し、
二重coneは元の対象の四つのshift・twistの直和となる。
従って元の対象は零である。一様な$r,s$は要求しない。
\end{proof}

\begin{lemma}[perfect性とSerre双対性からの正確な次元]
\label{CSBGlem:gluing-serre-dimension}
$A$は左右noetherian connected graded AS--Gorensteinで、
balanced dualizing complexが$A_\nu(-\ell)[d+1]$（$d\ge0$）で、
$\pi A\ne0$とする。
$\mathrm D^b(\operatorname{qgr}A)=\mathcal P_A$なら
$\operatorname{gldim}\operatorname{qgr}A=d$である。
\end{lemma}
\begin{proof}
\cite[Theorem~4.2, Corollary~4.3]{YZ}は左右の$\chi$と
有限cohomological dimensionを与える。
特に$\pi A(n)$は有限射影次元を持つ。
Chen\cite[Theorem~3.10]{CSBGChenGeneralizedSerre}の一般化Serre双対性は
有限射影次元の対象に適用でき、その関手は
$S(M)=M_\nu(-\ell)[d]$となる。
仮定により全coherent対象が適用範囲に入るので、
\[
 \operatorname{Ext}^i(M,N)^*
 \simeq\operatorname{Ext}^{d-i}(N,M_\nu(-\ell)).
\]
$i>d$では右辺が負次数ゆえ零。
また$M=\pi A$、$N=M_\nu(-\ell)$と取ると
$\operatorname{Ext}^d(M,N)^*\simeq\operatorname{Hom}(N,N)\ne0$なので等号となる。
\end{proof}

\begin{lemma}[超越的中心値の有限大域次元]
\label{CSBGlem:global-transcendental}
$k_0\subset\CC$、$R_0$は$k_0$-代数、$z\in Z(R_0)$は$k_0$上超越的とする。
$R=R_0\otimes_{k_0}\CC$が左右noetherianで$\operatorname{gldim}R\le n$なら、
$k_0$上超越的な全ての$\lambda\in\CC$について
$\operatorname{gldim}R/(z-\lambda)\le n$である。
\end{lemma}
\begin{proof}
$\lambda$を含む$\CC/k_0$の超越基底$T$を選び、$K=k_0(T\setminus\{\lambda\})$とする。
体拡大は自由かつ忠実平坦で、$K$-線形な$\CC\to K$の$1\mapsto1$という分裂がある。
従って大域次元の上界とnoetherian性は$R_K=R_0\otimes K$へ降下する。
$z$の冪は$K$上でも独立なので、中心局所化
$R_\eta=R_K\otimes_{K[z]}K(z)$も大域次元高々$n$。
$z\mapsto\lambda$で$K(z)=K(\lambda)$とみなすと
$R_\eta\otimes_{K(z)}\CC=R/(z-\lambda)$である。
最後の体拡大は代数的分離拡大で、右辺は$R$の商ゆえnoetherian。
\cite[Lemmas~3.2, 3.3]{RSVdB}の仮定が揃い、上界が保たれる。
任意の超越的体拡大が有限大域次元を保つと仮定したわけではない。
\end{proof}

\begin{theorem}[固定した良いambientでのvery-generalな正則性]
\label{CSBGthm:global-very-general}
$B$は中心一次元$u,v$を持つ左右noetherian AS正則五次元代数とし、
$G\in Z(B)_3$、$\operatorname{gr}_{x,y,z}B=S\otimes\CC[u,v]$とする。
$S$は整域、$\operatorname{gr}G=\gamma\in Z(S)_3$は非零正則元とする。
全ての$f\in\CC[u,v]_3$について$B/(G+f)$は$\CC[u,v]$上自由であり、
$\operatorname{qgr}(S/(\gamma))$は滑らかな楕円曲線のcoherent圏と仮定する。
有限表示の定義体$k_0$に対し、$a,d$が$k_0(b,c)$上代数的独立なら
\[
 \operatorname{gldim}\operatorname{qgr}
 B/(G+au^3+bu^2v+cuv^2+dv^3)=3.
\]
\end{theorem}
\begin{proof}
$D_u=B[u^{-1}]_0$はnoetherianで有限大域次元を持つ。
実際、中心局所化は大域次元を増やさず、
$B[u^{-1}]=D_u[u,u^{-1}]$から制限・誘導と評価$u=1$で上界を降下できる。
$t=v/u$とすると商の$u$-chartは
$D_u/(G/u^3+bt+ct^2+dt^3+a)$。
括弧内の$a$を除く中心元は$k_0(b,c,d)$上定義され、
最高主変数成分が非零正則なので、その冪は独立である。
前補題でchartの有限大域次元を得る。
$v$-chartも$d$を中心値として同様に扱える。
命題\ref{CSBGprop:global-two-charts}から有界 coherent 導来圏は twists で厚く生成される。
ここで $W=G+au^3+bu^2v+cuv^2+dv^3$ の最高成分は正則元 $\gamma$ なので、
$W$ 自身も正則である。正則中心商の Rees の補題
\cite[arXiv v1, Lemma~4.4(3)]{CSBGCCKM}により $A=B/(W)$ は
AS--Gorenstein 四次元であり、\cite[Corollary~4.3 とその直前]{YZ}から
balanced dualizing complex は $A_\nu(-\ell)[4]$ の形を持つ。
$A$ の $\CC[u,v]$-自由性から $\pi A\ne0$ でもある。
従って補題\ref{CSBGlem:gluing-serre-dimension}を $d=3$ として適用できる。
\end{proof}

\begin{lemma}[被約な係数環上の中心持上げ]
\label{CSBGlem:global-center-lifting}
$k$ を代数閉標数零体、$R_0$ を被約有限型 $k$-代数とし、
$U$ を非負の増大する exhaustive filtration $F_\bullet$ を持つ $R_0$-代数とする。
$R_0$ は中心に入り、$F_jU$ は有限自由、
$\operatorname{gr}U=H\otimes R_0$ が閉点への基底変更と両立すると仮定する。
$H=B(E,\mathcal L,\tau)$は滑らかな楕円曲線$E$，豊富な直線束$\mathcal L$，
有限位数の平行移動$\tau$によるtwisted homogeneous coordinate ringとする。
このとき $Z(H)$ の各斉次元は $Z(U)$ に持ち上がる。
さらに $U$ は有限型可換中心部分代数上有限である。
\end{lemma}

\begin{proof}
$H$ の有限個の斉次代数生成元を $U$ に持ち上げて $q_i$ と書く。
最高次 symbol を消す次数帰納法により、$q_i$ は $U$ を $R_0$-代数として生成する。
各整数 $j\geq 0$ に対し、有限自由 $R_0$-加群間の交換子写像
\[
 M_j:F_jU\longrightarrow\bigoplus_i F_{j+\deg q_i}U,
 \qquad w\longmapsto ([w,q_i])_i
\]
を取る。その核は $F_jU\cap Z(U)$ である。
Rains の定理 \cite[Theorem 2.13]{Rains}
は、標数零で有限位数の平行移動によりねじられた楕円代数の
\emph{任意の filtered deformation} の中心について
$\operatorname{gr}Z(U_s)=Z(H)$ を与える。
従って全ての閉点 $s$ で $\dim\ker(M_j\otimes k(s))$
は $\sum_{i\leq j}\dim Z(H)_i$ に等しい。
ここは単なる「associated graded が中心上有限」という一般論ではない。
楕円代数および有限位数平行移動という Rains の仮定を用いている。

$M_j$ の rank を $r_j$ とする。
$(r_j+1)$ 次 minor は全閉点で消える。
$R_0$ が被約かつ Jacobson なので、これらの minor は環の元として零である。
$r_j$ 次 minor の主開集合は $\operatorname{Spec}R_0$ を覆う。
各主開上で可逆な minor を用いて行列を変形すると、
$M_j$ は単位行列と零行列の直和になる。
従って $\ker M_j$ は有限射影で、その形成は任意の基底変更と可換する。

最高次 symbol は写像
\[
 \ker M_j\longrightarrow Z(H)_j\otimes R_0
\]
を与える。
この写像は全閉点で全射であるから、有限生成 cokernel に Nakayama を適用して全射である。
右辺は自由なので、各 $z\in Z(H)_j$ は $Z(U)\cap F_jU$ に持ち上がる。

最後に、有限位数 $r$ の平行移動に対し
$\mathcal L_r=\mathcal L\otimes\tau^*\mathcal L\otimes\cdots
\otimes(\tau^{r-1})^*\mathcal L$ は巡回置換の線型化を持ち、
$E/\langle\tau\rangle$ 上の豊富な直線束へ降下する。
その section ring は、次数を $r$ 倍にして $Z(H)$ に入る。
$H$ の各剰余次数部分は、有限射 $E\to E/\langle\tau\rangle$ の下で
coherent 層の section module と表されるので、$H$ はこの中心 section ring 上有限である。
従って $H$ が有限加群となる有限型斉次中心部分代数 $Z_0\subset Z(H)$ を取れる。
$Z_0$ の有限個の代数生成元と $H$ の有限個の $Z_0$-加群生成元を持ち上げる。
最高次 symbol を順次消して filtration 次数を下げる帰納法により、
持ち上げた中心元と $R_0$ が生成する可換部分代数上で $U$ は有限となる。
\end{proof}


\begin{theorem}[有限位数境界を持つ固定ambientの開な正則範囲]
\label{CSBGthm:global-open}
$k$ を代数閉標数零体とし、定理\ref{CSBGthm:global-very-general}の
良い固定ambientの仮定を $\CC$ の代わりに $k$ 上で満たす $B$ を取る。
さらに $H=S/(\gamma)=B(E,\mathcal L,\tau)$ のtranslationが有限位数とする。
ある非零多項式$P\in k[a,b,c,d]$が存在して、$P\ne0$上では
中心三次商のtails大域次元が3となる。
周囲のPBW/Koszul解消とnoetherian AS条件が基底変更と両立する構成では、
この十分条件は任意の代数閉拡大体上でも成立する。
特に $k=\CC$ で周囲の表示が有理数係数なら、有理数四係数の正則な商がZariski稠密に存在する。
\end{theorem}
\begin{proof}
$D_u=B[u^{-1}]_0$、$t=v/u$、$G_u=G/u^3$とする。
全係数を動かす$u$-chartの全代数は
\[
 \mathcal U=D_u[a,b,c,d]/(G_u+a+bt+ct^2+dt^3)
 \simeq D_u[b,c,d]
\]
なので左右noetherianかつ有限大域次元である。
$R=k[t,a,b,c,d]$をfiltration次数零とすると
$\operatorname{gr}\mathcal U=H\otimes R$で、各filtered pieceは有限自由。
前補題から$\mathcal U$は有限型可換中心部分代数$C_u$上有限である。
この$C_u$は係数環$R_0=k[a,b,c,d]$上も有限型である。
ここで\cite[Proposition~2.1, proof steps (6)--(8)]{RSVdB}を
$(R,S,A)=(R_0,C_u,\mathcal U)$へ適用する。
この命題は、有限型可換中心代数上の加群有限性と全代数の有限大域次元を仮定し、
$C_u$ の正規性や $\mathcal U$ が NCR であることを要求しない。
その証明は、ある非零 $P_u\in R_0$ を逆にした $R'=R_0[P_u^{-1}]$ 上で、
$\mathcal U'=\mathcal U\otimes_{R_0}R'$ の $R'$-平坦性と有限射影解消
\[
 0\longrightarrow Q_m\longrightarrow\cdots\longrightarrow Q_0
 \longrightarrow\mathcal U'\longrightarrow0,
 \qquad Q_i\in\operatorname{proj}
       ((\mathcal U')^{\mathrm{op}}\otimes_{R'}\mathcal U')
\]
を与える。
各 $Q_i$ と $\mathcal U'$ は $R'$-平坦であり、短完全列を順に見れば
全 syzygy も $R'$-平坦である。従って任意の体値 $R'\to L$ へのテンソル後も
これは対角 $\mathcal U'_L$ の有限射影解消である。
この解消は片側 $\mathcal U'_L$-加群として分裂するので、任意の
$\mathcal U'_L$-加群と片側でテンソルしても完全であり、
$L$ が体であることから長さ高々 $m$ の射影解消を得る。
ゆえに全ての幾何学的 fiber で有限大域次元が成立し、体拡大も許される。
$v$-chartでも$d$を消去して$P_v$を得る。
$P=P_uP_v$とし、二chartの貼合せとSerre双対性で結論する。
$k=\CC$ の場合、$\mathbb Q^4$は複素四次元affine空間でもZariski稠密なので、
非空開集合には有理数係数点が稠密に存在する。
この証明は Rains の標数零の中心持上げ、被約 Jacobson 環上の一定階数、
および有限中心代数に対する RSV の一般正則性だけを用い、
very-general 点の存在や $k$ の非可算性に依存しない。
反対環にも同じ議論を適用し、二つの非零多項式を掛ければ左右同時の結論を得る。
\end{proof}

\begin{lemma}[位数6の具体的な楕円曲線計算]\label{CSBGlem:order-six}
$E_1=V(x^3+y^3+z^3+xyz)$で$O=[1:-1:0]$を原点とすると、
$p=[1:-1:-1]$は位数6であり、$2p=[1:0:-1]$である。
\end{lemma}
\begin{proof}
$F=x^3+y^3+z^3+xyz$ と置く。
偏微分が全て零なら、いずれかの座標が零の場合は全座標が零となる。
全座標が非零の場合も、$3x^2=-yz$ とその巡回式を掛けると
$27x^2y^2z^2=-x^2y^2z^2$ となり矛盾する。従って $E_1$ は滑らかである。
$O$ の接線は $z=3(x+y)$ で、
\[
 F(x,y,3(x+y))=28(x+y)^3
\]
だから $O$ は変曲点である。
$p$ の接線 $2x+y+z=0$ については
\[
 F\bigl(-(y+z)/2,y,z\bigr)=\frac78(y+z)(y-z)^2.
\]
従って接線の第三交点は $r=[0:1:-1]$ であり、$p$ での交点重複度は正確に2である。
$r,O,q=[1:0:-1]$ は直線 $x+y+z=0$ 上にあるので、
楕円曲線の弦接線法から $2p=-r=q$ となる。
さらに $q$ の接線 $y=3(x+z)$ について
$F(x,3(x+z),z)=28(x+z)^3$ だから、$q\ne O$ は位数3の変曲点である。
従って $p$ の位数は3または6であるが、$p$ は変曲点でないので位数3ではなく、位数6となる。
\end{proof}
$h=-1$の場合は$x,y$の交換が$S_{-1}$と$S_1$の次数付き同型を与えるので、境界translationの位数も同じである。
